\chapter{Introduction}
\label{ch:intro}

\section{Motivation}
Tuberculosis is caused by the bacterium \emph{Mycobacterium tuberculosis}
and most often affects the lungs. According to the World Health
Organization, it was again the world's leading cause of death from a single
infectious agent in 2023, with an estimated 10.8 million people falling ill
and about 1.25 million deaths~\cite{who2024}. India carries the largest
share of the global burden, which makes early and cheap detection a
public-health priority in our own context.

Chest radiography is one of the main tools for TB screening. It is fast and
inexpensive, but every image must be read by a trained radiologist, and in
high-burden regions radiologists are scarce. The WHO therefore recommends
computer-aided detection software as an option for screening
adults~\cite{who2021screening}, and an independent evaluation of five
commercial products in Bangladesh found that all of them outperformed
experienced radiologists at the operating points studied~\cite{qin2021}.

Such systems are built on convolutional neural networks
(CNNs)~\cite{lecun1998}. Since a deep CNN won the ImageNet
challenge~\cite{krizhevsky2012, deng2009}, deep learning has become the
standard approach in computer vision~\cite{lecun2015} and in medical image
analysis~\cite{litjens2017}; large public collections such as
ChestX-ray8~\cite{wang2017} made it possible to train radiograph classifiers
that approach radiologist-level performance~\cite{rajpurkar2017}. For TB,
ImageNet-pretrained CNNs reach accuracies in the high nineties on public
datasets~\cite{rahman2020}.

Quantum computing offers a different model of computation. Feynman observed
that quantum systems are exponentially expensive to simulate
classically~\cite{feynman1982}; Shor's factoring algorithm~\cite{shor1997}
and Grover's search algorithm~\cite{grover1996} showed provable speed-ups
for specific problems; and in 2019 a 53-qubit processor performed a sampling
task argued to be beyond classical supercomputers~\cite{arute2019}. Present
devices, however, belong to the \emph{noisy intermediate-scale quantum}
(NISQ) era~\cite{preskill2018}: they have few qubits and every gate adds
noise. The dominant NISQ approach to machine learning is the
\emph{variational quantum circuit}, a short circuit with trainable rotation
angles optimised by a classical computer like the weights of a neural
network~\cite{schuld2020circuit, mitarai2018, farhi2018, benedetti2019}. The
same principle underlies the variational quantum
eigensolver~\cite{peruzzo2014, kandala2017}; the family of methods is
reviewed in~\cite{cerezo2021vqa, schuld2018book}.

Because current devices are small, practical models are \emph{hybrid}: a
classical network handles the high-dimensional image and a small circuit is
inserted at one point. In \emph{classical-to-quantum} transfer learning a
CNN compresses the image into a few features that a circuit
classifies~\cite{mari2020}; in a \emph{quanvolutional} network a circuit is
slid over the image like a convolution filter~\cite{henderson2020}.
PennyLane~\cite{bergholm2018} makes both kinds of model ordinary
PyTorch~\cite{paszke2019} modules that can be trained by backpropagation
through a simulator.

\section{Problem Statement}
Published hybrid models are often compared only with a classical model of
a different size or structure, on a single run, so it is hard to tell
whether the quantum circuit is responsible for the result. This project
asks:

\begin{quote}
\emph{Does adding a small quantum circuit improve a CNN classifier for
tuberculosis detection from chest X-rays, and does the answer depend on
whether the circuit is placed after the classical feature extractor (\cq)
or before it (\qc), once each hybrid is compared with a classical model of
the same structure and size?}
\end{quote}

\section{Objectives and Contributions}
The objectives were to build a reproducible pipeline for the TB Chest X-ray
Database; to implement a classical CNN baseline, a \cq{} hybrid and a \qc{}
hybrid together with classical controls matched to each hybrid; to study
the effect of circuit design, backbone and training-set size with repeated
seeds and paired statistical tests; and to analyse the circuits themselves.
The main contributions are:
\begin{enumerate}
  \item A configurable PyTorch--PennyLane code base in which every
        experiment differs from its reference in exactly one setting.
  \item A comparison of both hybrid placements with a baseline and four
        matched controls on five training-set sizes and three seeds
        (105 training runs), a 12-configuration circuit ablation, and a
        backbone $\times$ head study including pretrained ResNet18.
  \item A quantitative analysis of the circuits: expressibility,
        entangling capability, barren-plateau scaling, and robustness to
        finite shots and depolarising noise.
  \item An honest assessment: the \cq{} hybrid is consistently but not
        significantly better than its controls, and the \qc{} hybrid's
        parameter efficiency is reproduced by classical downsampling.
\end{enumerate}

Chapter~\ref{ch:related} reviews the background and related work,
Chapter~\ref{ch:proposed} describes the data, models and protocol,
Chapter~\ref{ch:results} presents and analyses the results, and
Chapter~\ref{ch:conclusion} concludes.
